% Instalacja MESA i minimalny przyklad ewolucji gwiazdy o masie 1 Msun
% Linux, macOS, Windows (WSL), Windows natywnie, pracownia komputerowa.
% Kompilacja: pdflatex MESA_instrukcja.tex (dwukrotnie, dla spisu tresci)
\documentclass[11pt,a4paper]{article}

\usepackage[T1]{fontenc}
\usepackage[utf8]{inputenc}
\usepackage[polish]{babel}
\usepackage{lmodern}
\usepackage{textcomp}
\usepackage{upquote}
\usepackage{microtype}
\usepackage[a4paper,margin=2.2cm]{geometry}
\usepackage{xcolor}
\usepackage{listings}
\usepackage{enumitem}
\usepackage{booktabs}
\usepackage{array}
\usepackage{longtable}
\usepackage{tcolorbox}
\usepackage[hidelinks]{hyperref}

\setlength{\parindent}{0pt}
\setlength{\parskip}{0.55em}
\setlength{\emergencystretch}{3em}

% ---------------------------------------------------------------- notatki robocze
% \uwagitrue  -> czerwone notatki "do sprawdzenia" widoczne (wersja robocza)
% \uwagifalse -> notatki ukryte (wersja dla studentow)
\newif\ifuwagi
\uwagitrue
\newcommand{\dosprawdzenia}[1]{\ifuwagi\textcolor{red}{\small[do sprawdzenia: #1]}\fi}

\newcommand{\repourl}{\url{https://github.com/VA00/MESA-minimal-test-1MSun-star}}

% ---------------------------------------------------------------- listingi
\definecolor{codebg}{gray}{0.95}
\definecolor{coderule}{gray}{0.70}
\lstdefinestyle{cmd}{
  basicstyle=\ttfamily\small, columns=fullflexible, keepspaces=true, upquote=true,
  showstringspaces=false, breaklines=false, frame=single, framerule=0.4pt,
  rulecolor=\color{coderule}, backgroundcolor=\color{codebg},
  xleftmargin=3pt, xrightmargin=3pt, framesep=4pt, aboveskip=0.6em, belowskip=0.6em}
\lstdefinestyle{out}{
  basicstyle=\ttfamily\small, columns=fullflexible, keepspaces=true, upquote=true,
  showstringspaces=false, breaklines=false, frame=leftline, framerule=1.5pt,
  rulecolor=\color{coderule}, xleftmargin=8pt, framesep=6pt, aboveskip=0.6em, belowskip=0.6em}
% szeroki blok (wysuniety na marginesy) dla dlugich polecen jednolinijkowych
\lstdefinestyle{cmdwide}{style=cmd, basicstyle=\ttfamily\footnotesize,
  xleftmargin=-14mm, xrightmargin=-14mm}
\lstnewenvironment{shw}{\lstset{style=cmdwide}}{}
\lstnewenvironment{psw}{\lstset{style=cmdwide}}{}
\lstnewenvironment{sh}{\lstset{style=cmd}}{}   % bash / zsh
\lstnewenvironment{ps}{\lstset{style=cmd}}{}   % PowerShell 7
\lstnewenvironment{out}{\lstset{style=out}}{}  % wydruk, zawartosc plikow
\newcommand{\code}[1]{\texttt{#1}}

\newtcolorbox{uwaga}[1][Uwaga]{colback=yellow!6, colframe=orange!70!black, fonttitle=\bfseries,
  title=#1, boxrule=0.5pt, arc=1pt, left=4pt, right=4pt, top=2pt, bottom=2pt}
\newtcolorbox{info}[1][Informacja]{colback=blue!3, colframe=blue!50!black, fonttitle=\bfseries,
  title=#1, boxrule=0.5pt, arc=1pt, left=4pt, right=4pt, top=2pt, bottom=2pt}

\setlist{itemsep=0.2em, topsep=0.3em}

\title{Instalacja i minimalny przykład komputerowej ewolucji gwiazdy w~programie MESA\\[0.3em]
  \Large Linux, macOS, Windows (WSL i natywnie)}
\author{Andrzej Odrzywołek\\ \small\href{mailto:andrzej.odrzywolek@uj.edu.pl}{andrzej.odrzywolek@uj.edu.pl}}
\date{wersja robocza, 7~października 2026}

\begin{document}
\maketitle

\begin{info}[O tej instrukcji]
Instrukcja i wszystkie potrzebne pliki znajdują się w repozytorium\\
\repourl\\
Przykład to absolutnie minimalna konfiguracja MESA r24.08.1: ewolucja gwiazdy o masie
$1\,M_\odot$ od jednorodnego obłoku H/He (pre-MS) do stygnącego białego karła, ok.~14\,000 kroków.
Obliczenia trwają od ok.~1~godziny (szybki komputer) do kilku godzin.

\textbf{Którą część czytać?}
\begin{itemize}[leftmargin=1.2em]
  \item Pracownia komputerowa Fort OA (MESA już zainstalowana): rozdział~\ref{sec:pracownia}.
  \item Własny komputer z Linuksem: rozdział~\ref{sec:linux}.
  \item Mac z procesorem Apple Silicon (M1--M4): rozdział~\ref{sec:mac}.
  \item Windows: \textbf{zalecane} WSL (pełna funkcjonalność, z wykresami), rozdział~\ref{sec:wsl}.
        Bez WSL: wersja natywna (eksperymentalna, na razie bez wykresów), rozdział~\ref{sec:win}.
\end{itemize}
Potem: uruchomienie przykładu (rozdział~\ref{sec:run}), animacje (\ref{sec:anim}) i tabela
benchmarków (\ref{sec:bench}).
\end{info}

\tableofcontents

% =====================================================================
\section{Wymagania}\label{sec:wymagania}

\begin{itemize}
  \item Dysk: pełna instalacja MESA to ok.~21~GB, MESA SDK ok.~0,7~GB.
  \item Pamięć: minimum teoretyczne 8~GB RAM, w praktyce nieco więcej.
  \item Ścieżki do katalogów z MESA \textbf{nie mogą zawierać spacji}.
\end{itemize}

% =====================================================================
\section{Pracownia komputerowa Fort OA}\label{sec:pracownia}

MESA jest zainstalowana systemowo w~\verb!/opt/MESA!. Wystarczającą ilość pamięci mają w~tej
chwili komputery: Cep, Aqr, Eri, Sgr, Hya.

\begin{enumerate}[leftmargin=1.5em, itemsep=0.6em]
\item Zmienne środowiskowe dla MESA SDK i programu MESA:
\begin{sh}
export MESASDK_ROOT=/opt/MESA/mesasdk
source $MESASDK_ROOT/bin/mesasdk_init.sh
export MESA_DIR=/opt/MESA/mesa-24.08.1
\end{sh}
\item Aby uniknąć problemów z zapisem plików cache (eos -- równanie stanu, kap --
  nieprzezroczystość, rates -- reakcje jądrowe, ionization -- jonizacja), tworzymy własny katalog
  i ustawiamy zmienną \code{MESA\_CACHES\_DIR}:
\begin{sh}
mkdir -p $HOME/.mesa
export MESA_CACHES_DIR=$HOME/.mesa
\end{sh}
\item Pracujemy w katalogu \verb!/data/rok4! (zamień \code{rok4} na swój rocznik). Testy pokazują,
  że symulacja działa tam ok.~3 razy szybciej niż w~katalogu \verb!$HOME!, ale wyniki nie są
  dostępne na innych komputerach.
\begin{sh}
cd /data/rok4
mkdir MESA2027
cd MESA2027
git clone https://github.com/VA00/MESA-minimal-test-1MSun-star.git
cd MESA-minimal-test-1MSun-star
\end{sh}
\item Kompilacja i uruchomienie: \verb!./mk!, \verb!./rn! (rozdział~\ref{sec:run}).
\end{enumerate}

\begin{uwaga}[Znany problem]
Na komputerach pracowni obliczenia zawieszały się przy kroku 10\,703 (okolice błysku helowego).
\end{uwaga}

% =====================================================================
\section{Linux (Ubuntu) -- instalacja na własnym komputerze}\label{sec:linux}

Aktualizacja i sprawdzenie: 18 października 2025 (MESA SDK 24.7.1, MESA 24.08.1).

\begin{enumerate}[leftmargin=1.5em, itemsep=0.6em]
\item Pobieramy MESA SDK (24.7.1, ok.~209~MB, ok.~8~min) i MESA (24.08.1, ok.~2~GB, ok.~18~min).
  Serwer SDK odrzuca pobieranie z konsoli, dlatego \code{wget} „udaje” przeglądarkę Firefox:
\begin{shw}
mkdir ~/MESA
cd ~/MESA
wget -U firefox http://user.astro.wisc.edu/~townsend/resource/download/mesasdk/mesasdk-x86_64-linux-24.7.1.tar.gz
wget https://zenodo.org/records/13353788/files/mesa-24.08.1.zip
\end{shw}
\item Rozpakowujemy SDK i ustawiamy zmienne (bash):
\begin{sh}
tar xvfz mesasdk-x86_64-linux-24.7.1.tar.gz -C ~/MESA/
export MESASDK_ROOT=~/MESA/mesasdk
source $MESASDK_ROOT/bin/mesasdk_init.sh
\end{sh}
  Powinny pojawić się komunikaty:
\begin{out}
mesasdk_init.sh: checking architecture
mesasdk_init.sh: regenerating headers
mesasdk_init.sh: checking prerequisites
\end{out}
\item Na ogół brakuje kilku programów, o czym informuje komunikat typu:
\begin{out}
mesasdk_init.sh: missing prerequisites:
  C shell Z library X11 library
\end{out}
  Doinstalowujemy je i powtarzamy krok~2 (\code{source ...}):
\begin{sh}
sudo apt install tcsh zlib1g-dev libx11-dev
\end{sh}
\item Rozpakowujemy MESA i ustawiamy kolejne zmienne (liczba wątków według liczby rdzeni):
\begin{sh}
cd ~/MESA
unzip mesa-24.08.1.zip
export MESA_DIR=~/MESA/mesa-24.08.1
export OMP_NUM_THREADS=4
\end{sh}
\item Opcjonalnie, żeby zmienne ustawiały się przy każdym uruchomieniu terminala, dopisujemy je do
  \verb!~/.bashrc!:
\begin{sh}
echo 'export MESASDK_ROOT=~/MESA/mesasdk' >> ~/.bashrc
echo 'source $MESASDK_ROOT/bin/mesasdk_init.sh' >> ~/.bashrc
echo 'export MESA_DIR=~/MESA/mesa-24.08.1' >> ~/.bashrc
echo 'export OMP_NUM_THREADS=4' >> ~/.bashrc
\end{sh}
\item Instalacja (kompilacja trwa kilkanaście minut):
\begin{sh}
cd $MESA_DIR
./install
\end{sh}
  Jeżeli się powiodła, pojawi się dużo gwiazdek i napis:
\begin{out}
************************************************
MESA installation was successful
************************************************
\end{out}
\item Pobieramy przykład:
\begin{sh}
cd ~/MESA
git clone https://github.com/VA00/MESA-minimal-test-1MSun-star.git
cd MESA-minimal-test-1MSun-star
\end{sh}
  i uruchamiamy go poleceniami \verb!./mk!, \verb!./rn! (rozdział~\ref{sec:run}).
\end{enumerate}

\begin{info}[Instalacja na dysku zewnętrznym]
Na systemach plików bez dowiązań symbolicznych (exFAT, FAT32) SDK i~MESA nie działają. Można je
zainstalować na dysku lokalnym, a potem skopiować na dysk USB z~rozwinięciem dowiązań:
\begin{sh}
rsync -aL --info=progress2 ~/MESA/mesa-24.08.1/ /media/$USER/USB/mesa-24.08.1/
rsync -aL --info=progress2 ~/MESA/mesasdk/ /media/$USER/USB/mesasdk/
\end{sh}
i ustawić \code{MESA\_DIR} oraz \code{MESASDK\_ROOT} na nowe położenie.
\end{info}

% =====================================================================
\section{macOS (Apple Silicon)}\label{sec:mac}

Aktualizacja: 12 października 2025 (MESA SDK 25.8.1, MESA 24.08.1).

\begin{enumerate}[leftmargin=1.5em, itemsep=0.6em]
\item Pobieramy i instalujemy MESA SDK dla Apple Silicon, wersja 25.8.1 (sierpień 2025), z katalogu\\
  {\small\url{http://user.astro.wisc.edu/~townsend/resource/download/mesasdk/}}\\
  plik \code{mesasdk-aarch64-macos-25.8.1.pkg}.
  Instalacja jest okienkowa, po kliknięciu na pobraną „paczkę”. SDK trafia do
  \verb!/Applications/mesasdk!. Starsza wersja 24.10.1 już nie działa.
\item Zmienne środowiskowe (domyślna powłoka zsh jest zgodna z bash):
\begin{sh}
export MESASDK_ROOT=/Applications/mesasdk
source $MESASDK_ROOT/bin/mesasdk_init.sh
\end{sh}
\item Pobieramy MESA (plik zip z Zenodo, \url{https://zenodo.org/records/13353788}). Safari sam
  rozpakowuje archiwum. Cały katalog \code{mesa-24.08.1} najlepiej przenieść do \verb!~/MESA!.
  To samo z terminala:
\begin{sh}
mkdir ~/MESA
cd ~/MESA
curl -L -O https://zenodo.org/records/13353788/files/mesa-24.08.1.zip
unzip mesa-24.08.1.zip
\end{sh}
\item Kolejne zmienne:
\begin{sh}
export MESA_DIR=$HOME/MESA/mesa-24.08.1
export OMP_NUM_THREADS=8
export PATH=$PATH:$MESA_DIR/scripts/shmesa
\end{sh}
  Żeby ustawiały się automatycznie, dopisujemy powyższe linie (także te z kroku~2) do
  \verb!~/.zshrc!.
\item Instalacja:
\begin{sh}
cd $MESA_DIR
./install
\end{sh}
\item Aby MESA cokolwiek wyświetlała (okna pgstar), trzeba zainstalować serwer X, np.~XQuartz
  (\url{https://www.xquartz.org}).
\item Pobieramy przykład i uruchamiamy go jak w Linuksie:
\begin{sh}
cd ~/MESA
git clone https://github.com/VA00/MESA-minimal-test-1MSun-star.git
cd MESA-minimal-test-1MSun-star
\end{sh}
  Czas obliczeń: 135~min przy \code{OMP\_NUM\_THREADS=8} na procesorze M3 Max.
\end{enumerate}

\begin{info}[Do zbadania]
Nowoczesne procesory (Apple M3, M4, ale też np.~Intel i9 285) mają rdzenie typu
„performance” (szybkie, gorące) i~„efficiency” (wolniejsze, chłodniejsze). W Linuksie można
przypisać symulację do wybranych rdzeni poleceniem \code{taskset}. Jak zrobić to w macOS?
\end{info}

% =====================================================================
\section{Windows -- WSL (zalecane)}\label{sec:wsl}

\dosprawdzenia{rozdział przygotowany na podstawie dokumentacji WSL, na komputerze testowym nie było WSL}

\begin{enumerate}[leftmargin=1.5em, itemsep=0.6em]
\item W PowerShell \textbf{uruchomionym jako administrator}, potem restart komputera:
\begin{ps}
wsl --install
\end{ps}
  W oknie Ubuntu, które otworzy się po restarcie, zakładamy użytkownika i hasło.
\item Pamięć dla WSL. W zwykłym PowerShell:
\begin{ps}
notepad $env:USERPROFILE\.wslconfig
\end{ps}
  wpisujemy (np.~dla komputera z 16~GB RAM)
\begin{out}
[wsl2]
memory=12GB
\end{out}
  zapisujemy plik i wykonujemy \verb!wsl --shutdown!.
\item W oknie Ubuntu:
\begin{sh}
sudo apt update
sudo apt install tcsh zlib1g-dev libx11-dev binutils make perl unzip wget
\end{sh}
  i dalej według rozdziału~\ref{sec:linux}, w katalogu \verb!~/MESA! (\textbf{nie} w \verb!/mnt/c!).
\end{enumerate}

% =====================================================================
\section{Windows -- natywnie (PowerShell 7, bez WSL)}\label{sec:win}

Wersja natywna nie wyświetla wykresów (pgstar). Wszystkie polecenia wpisujemy w~PowerShell~7.

\begin{enumerate}[leftmargin=1.5em, itemsep=0.6em]
\item PowerShell~7:
\begin{ps}
winget install --id Microsoft.PowerShell -e
\end{ps}
  Dalej pracujemy w oknie „PowerShell 7”.
\item MSYS2 i kompilator gfortran:
\begin{psw}
winget install --id MSYS2.MSYS2 -e
C:\msys64\usr\bin\pacman.exe -Syu --noconfirm
C:\msys64\usr\bin\pacman.exe -Syu --noconfirm
C:\msys64\usr\bin\pacman.exe -S --needed --noconfirm mingw-w64-ucrt-x86_64-gcc-fortran
C:\msys64\usr\bin\pacman.exe -S --needed --noconfirm mingw-w64-ucrt-x86_64-make
C:\msys64\usr\bin\pacman.exe -S --needed --noconfirm mingw-w64-ucrt-x86_64-hdf5
\end{psw}
\item Python i Git, a potem \textbf{nowe okno} PowerShell~7:
\begin{ps}
winget install --id Python.Python.3.14 -e
winget install --id Git.Git -e
\end{ps}
\item Pobranie i rozpakowanie MESA (komunikaty \verb!Can't create ...! można zignorować):
\begin{ps}
mkdir C:\MESA
cd C:\MESA
curl.exe -L -O https://zenodo.org/records/13353788/files/mesa-24.08.1.zip
tar -xf mesa-24.08.1.zip
\end{ps}
  W sieci OA zamiast Zenodo:
\begin{ps}
curl.exe -L -O http://wolarz.oa.uj.edu.pl/repo/mesa-24.08.1.zip
\end{ps}
\item Pobranie przykładu:
\begin{ps}
git clone https://github.com/VA00/MESA-minimal-test-1MSun-star.git
\end{ps}
\item Zmienne środowiskowe, \textbf{w każdym nowym oknie} PowerShell:
\begin{ps}
. C:\MESA\MESA-minimal-test-1MSun-star\windows\mesa_env.ps1
\end{ps}
\item Kompilacja MESA (kilka minut, na końcu napis \code{MESA installation was successful}):
\begin{ps}
C:\MESA\MESA-minimal-test-1MSun-star\windows\install_mesa.ps1
\end{ps}
\item Kompilacja i uruchomienie przykładu:
\begin{ps}
cd C:\MESA\MESA-minimal-test-1MSun-star
./mk
./rn
\end{ps}
\end{enumerate}
Dalej jak w rozdziale~\ref{sec:run}.

% =====================================================================
\section{Uruchomienie przykładu (wszystkie systemy)}\label{sec:run}

\subsection{Polecenia}

\begin{center}
\begin{tabular}{@{}ll@{}}
\toprule
kompilacja & \verb!./mk! \\
start od początku & \verb!./rn! \\
restart od ostatniego zdjęcia & \verb!./re! \\
restart od zdjęcia \code{x500} & \verb!./re x500! \\
sprzątanie kompilacji & \verb!./clean! \\
\bottomrule
\end{tabular}
\end{center}
Obliczenia przerywamy kombinacją \code{Ctrl+C}. Zdjęcia do restartów są w katalogu \code{photos}.

\subsection{Szybki test: 50 kroków od ciągu głównego}\label{sec:zams}

Po kompilacji (\verb!./mk!):
\begin{sh}
mkdir ../zams_test
cp star inlist inlist.pgstar history_columns.list rn ../zams_test/
cd ../zams_test
\end{sh}
W Windows natywnie:
\begin{ps}
mkdir ../zams_test
Copy-Item star.exe, inlist, inlist.pgstar, history_columns.list, rn.ps1 ../zams_test
cd ../zams_test
\end{ps}
W pliku \code{inlist} (edytor \code{nano inlist} lub \code{notepad inlist}) zmieniamy w sekcji
\code{\&star\_job}
\begin{out}
  create_pre_main_sequence_model = .false.
\end{out}
a w sekcji \code{\&controls} dopisujemy linijkę
\begin{out}
  max_model_number = 50
\end{out}
Uruchamiamy \verb!./rn!. Poprawny wynik (0 powtórzeń, 50 kroków):
\begin{out}
termination code: max_model_number
                 runtime (minutes), retries, steps        0.10         0        50
\end{out}
\dosprawdzenia{wynik sprawdzony w Windows natywnie; w Linuksie/macOS spodziewany taki sam}
\subsection{Pełna ewolucja}

Po \verb!./rn! po kilku minutach pojawiają się okna z wizualizacją ewolucji gwiazdy (poza wersją
natywną Windows). Koniec obliczeń:
\begin{out}
save photos/x122 for model 14122
termination code: log_L_lower_limit

                 runtime (minutes), retries, steps      113.25       303     14122
\end{out}
(MacBook M3 Max, 10~wątków). Wyniki trafiają do katalogów \code{LOGS} (\code{history.data},
profile), \code{photos} (zdjęcia do restartów) i \code{png} (wykresy).

W trakcie ewolucji możemy edytować plik \code{inlist.pgstar}, w szczególności dodawać różne sposoby
wizualizacji.

\begin{uwaga}[Znane problemy]
\begin{itemize}[leftmargin=1.2em]
  \item Błąd \code{Cannot match namelist object name read\_extra\_pgstar\_inlist1}: w pliku
        \code{inlist} trzeba zmienić \code{read\_extra\_pgstar\_inlist1} na
        \code{read\_extra\_pgstar\_inlist(1)} oraz \code{extra\_pgstar\_inlist1\_name} na
        \code{extra\_pgstar\_inlist\_name(1)} (w repozytorium już poprawione).
  \item Błąd związany z \code{adipls} może wynikać ze zmiany nazewnictwa nowych wersji MESA: nie
        ma literki „r”, tylko sam numer wersji.
\end{itemize}
\end{uwaga}

% =====================================================================
\section{Animacje}\label{sec:anim}

Z plików w katalogu \code{png} można utworzyć film programem \code{ffmpeg}:
\begin{sh}
# film AVI
ffmpeg -f image2 -pattern_type glob -i 'png/*.png' -c:v libx264 -preset slow \
  -crf 20 -pix_fmt yuv420p -aspect 16:9 -vf "scale=1920:-1" -r 25 movie.avi

# film MKV bezstratny
ffmpeg -f image2 -pattern_type glob -i 'png/*.png' -c:v ffv1 -level 3 \
  -pix_fmt yuv444p -aspect 16:9 -vf "scale=1200:-1" -r 25 movie.mkv

# animacja GIF
ffmpeg -y -f image2 -thread_queue_size 256 -pattern_type glob -i 'png/*.png' \
  -vf "scale=800:-1:flags=lanczos,palettegen" /tmp/palette.png
ffmpeg -y -f image2 -thread_queue_size 256 -pattern_type glob -i 'png/*.png' \
  -i /tmp/palette.png \
  -lavfi "scale=1600:-1:flags=lanczos [x]; [x][1:v] paletteuse" movies.gif
\end{sh}
Skrypt \verb!./parallel_movie_encode.sh! tworzy równolegle osobne animacje dla każdego typu
wizualizacji.

% =====================================================================
\section{Tabela benchmarków}\label{sec:bench}

Czas obliczeń testowego modelu (w minutach) na różnych komputerach i przy różnej liczbie wątków
\code{OMP\_NUM\_THREADS}. Czas jest podany w końcowych linijkach wydruku:
\code{runtime (minutes), retries, steps}.

{\small
\begin{longtable}{@{}>{\raggedright\arraybackslash}p{0.36\textwidth}>{\raggedright\arraybackslash}p{0.58\textwidth}@{}}
\toprule
\textbf{Komputer} & \textbf{wątki: czas [min]} \\
\midrule
\endhead
MacBook M3 Max & 1: 747,52; 2: 391,84; 5: 181,67; 6: 160,08; 7: 144,87; 8: 134,33; 9: 123,51; 10: 113,25 \\
VMware & 4: 157,74 \\
Pracownia OA Fort (Cep, Aqr, Eri, Sgr, Hya) & 4: zawiesił się \\
Intel Core Ultra 9 285K, 128 GB RAM, Windows 11, WSL2 & 1: 411,39; 2: 220,32; 3: 155,52; 4: 120,57; 5: 99,37; 6: 85,45; 8: 70,12; 24: 42,39 \\
Super Micro, 2× Intel Xeon Gold 6238R 2,2 GHz (2×26 rdzeni), 376 GB RAM, Ubuntu 22.04 & 1: 732,24; 2: 397,55; 3: 283,02; 4: 223,56; 5: 184,39; 6: 161,15; 7: 143,23; 8: 129,57; 9: 118,77; 10: 109,70; 24: 67,23 \\
AMD Ryzen 7 7435HS 3,1 GHz, 16 GB RAM, Windows 11, WSL2 & 2: 393,79; 4: 198,33; 8: 152,86 \\
MacBook Air M4, 16 GB RAM & 2: nie skończył przez całe zajęcia \\
Intel i7-4770K 3,4 GHz, 16 GB RAM, Ubuntu & 1: 783,99; 2: 420,83; 3: 308,45; 4: 246,79; 5: 237,08; 6: 227,10; 7: 218,53; 8: 210,82 \\
AMD Ryzen 9 5900X, 32 GB RAM, Windows 11 natywnie (gfortran 16.1 MinGW) & 12: 70,77 \\
\bottomrule
\end{longtable}
}

% =====================================================================
\appendix
\section{Windows natywnie: szczegóły i rozwiązywanie problemów}\label{app:win}

Sprawdzone 6--7~października 2026: Windows~11, AMD Ryzen~9 5900X, 32~GB RAM, gfortran 16.1.0
(MSYS2), HDF5 2.1.1. Ograniczenia: brak wykresów pgstar (MESA zbudowana bez PGPLOT,
\code{pgstar\_flag} jest ignorowane), nie są budowane GYRE, ADIPLS, \code{binary} i \code{astero},
wyniki nie są identyczne bit w bit z~Linuksem (różnice na poziomie zaokrągleń).

\subsection{Testy modułów (opcjonalnie)}

Po instalacji MESA, w oknie z ustawionymi zmiennymi (kilkanaście minut):
\begin{ps}
C:\MESA\MESA-minimal-test-1MSun-star\windows\test_mesa.ps1
\end{ps}
Wyniki są porównywane z wzorcami z Linuksa (tolerancja $10^{-8}$). Oczekiwane: \code{OK} wszędzie
oprócz \code{num} (całkowanie sztywnych równań różniczkowych wzmacnia różnice zaokrągleń do
$10^{-5}$) i \code{net} (jedna linia, różnica $5\cdot10^{-7}$).

\subsection{Parametry \texttt{mesa\_env.ps1}}

\code{-MesaDir} (inne położenie MESA niż katalog obok repozytorium), \code{-Threads}
(liczba wątków, domyślnie liczba rdzeni), \code{-Toolchain} (MSYS2 w innym miejscu niż
\verb!C:\msys64!).

\subsection{Co zmieniono względem instalacji na Linuksie}

\begin{enumerate}[leftmargin=1.5em, itemsep=0.4em]
\item \textbf{System budowania.} Moduły są kompilowane oryginalnymi makefile'ami MESA, uruchamianymi
  natywnym \code{mingw32-make.exe} (receptury wykonuje \code{cmd.exe}), z plikiem
  \verb!utils\makefile_header! w wersji dla Windows. Narzędzia z~MESA SDK zastąpiono skryptami
  PowerShell (\code{makedepf90.ps1}, \code{cp\_if\_newer.ps1}), a bashowe \code{install}
  i~\code{build\_data\_and\_export} skryptem \code{install\_mesa.ps1}. Biblioteka forum (z GYRE)
  jest preprocesowana programem \code{fypp}. LAPACK/BLAS pochodzą ze źródeł w~\code{mesa/mtx},
  HDF5 z~MSYS2 (linkowana statycznie). GYRE jest zastąpione stubem.
\item \textbf{Plik C} \verb!utils\private\utils_c_system.c! używa funkcji POSIX, których MinGW nie
  ma w tej postaci (\code{mkdir(path, mode)}, \code{realpath}, brak \code{O\_BINARY}). Łata jest
  ujęta w \verb!#ifdef _WIN32!.
\item \textbf{Błąd ustawień regionalnych w libgfortran (MinGW/UCRT).} Wokół każdej instrukcji
  wejścia-wyjścia libgfortran przełącza locale na „C”, po czym „przywraca” poprzednie wskaźnikiem
  do już zwolnionej pamięci. Proces przechodzi wtedy na locale użytkownika
  (\code{Polish\_Poland.1250}, przecinek dziesiętny). W wątkach OpenMP liczba „0.6000” jest
  odczytywana jako 0 i MESA przerywa wczytywanie tablic EOS (\code{bad header info}). Na Windows
  z angielskimi ustawieniami regionalnymi błąd się nie ujawnia. Poprawka:
  \code{mesa\_win\_locale.c} przechwytuje \code{setlocale}
  (\verb!-Wl,--wrap=__imp_setlocale!, libgfortran statycznie) i utrzymuje locale „C”. Przy okazji
  odczyt plików tekstowych jest ok.~2 razy szybszy. \code{libgomp} i \code{libwinpthread} pozostają
  dynamiczne, bo wersja w pełni statyczna czytała pliki ok.~2,5 razy wolniej.
\item \textbf{Stos.} Windows daje głównemu wątkowi 2~MB stosu (Linux: 8~MB lub
  \code{ulimit -s unlimited}). \code{star.exe} jest linkowany z rezerwacją 1~GB
  (\verb!-Wl,--stack,1073741824!), a~wątkom OpenMP \code{mesa\_env.ps1} ustawia
  \code{OMP\_STACKSIZE=512M}.
\item \textbf{Błąd optymalizacji gfortrana 16.1.} Plik \verb!star\private\star_bcyclic.f90!
  (blokowo-trójdiagonalny solver równań budowy gwiazdy, z dołączanym \code{mtx\_solve\_routines.inc})
  skompilowany z \code{-O1}/\code{-O2} daje błędne wyniki: już pierwsza iteracja Newtona jest
  rozbieżna, nawet przy starcie z modelu ZAMS. Plik wskazała bisekcja po 159 obiektach modułu
  \code{star}. Jest kompilowany z \code{-Og} (poziom optymalizacji „pod debugowanie”, między
  \code{-O0} a~\code{-O1}), reszta MESA z~\code{-O2}. Przyczyna (błąd GCC czy niezdefiniowane
  zachowanie w~kodzie) nie jest jeszcze ustalona.
\item \textbf{Drobiazgi.} W \verb!chem\make\makefile_base! usunięto \verb!-m %m.mod! (\code{cmd.exe}
  traktuje \verb!%m! jak zmienną). Dowiązanie \code{pgstar.f90} zastępuje kopia
  \code{pgstar\_stub.f90}. \code{star.exe} potrzebuje \code{libgomp-1.dll}
  i~\code{libwinpthread-1.dll} z~\verb!C:\msys64\ucrt64\bin! (dodaje je do \code{PATH}
  \code{mesa\_env.ps1}).
\end{enumerate}

\subsection{Rozwiązywanie problemów}

\begin{description}[style=nextline, leftmargin=1.2em, itemsep=0.3em]
\item[\code{MESA\_DIR is not set}, \code{mingw32-make} nie jest rozpoznawane, okno „nie znaleziono
      libgomp-1.dll”]
  W tym oknie nie wykonano \verb!. ...\windows\mesa_env.ps1! (z kropką na początku). MSYS2 w~innym
  miejscu niż \verb!C:\msys64!: parametr \verb!-Toolchain D:/msys64/ucrt64!.
\item[Skrypty są blokowane („running scripts is disabled on this system”)]
\begin{ps}
Set-ExecutionPolicy -Scope CurrentUser RemoteSigned
\end{ps}
\item[\code{pacman}: błędy pobierania pakietów (404, \code{failed retrieving file})]
  Nieaktualna baza pakietów: \verb!C:\msys64\usr\bin\pacman.exe -Syu --noconfirm! i ponowna
  instalacja.
\item[\code{tar}: \code{Error opening archive} przy pliku \code{.zip}]
  Przed systemowym \code{tar} w \code{PATH} jest GNU tar (z Git lub MSYS2). Pełna ścieżka:
  \verb!C:\Windows\System32\tar.exe -xf mesa-24.08.1.zip!.
\item[\code{Python 3 not found}, otwiera się Microsoft Store]
  Zainstaluj Pythona i otwórz nowe okno PowerShell. Alias Sklepu można wyłączyć w~\emph{Ustawienia
  → Aplikacje → Zaawansowane ustawienia aplikacji → Aliasy wykonywania aplikacji}.
\item[Rozbieżność już w pierwszym kroku (\code{correction rejected by sizeB})]
  Kończy się komunikatem \code{dt < min\_timestep\_limit}. Objaw błędu z punktu~5, np.~gdy MESA
  była skompilowana ze starszą wersją skryptów.
  Przebuduj moduł \code{star} i katalog roboczy:
\begin{ps}
C:\MESA\MESA-minimal-test-1MSun-star\windows\install_mesa.ps1 -Modules star -Clean
./clean
./mk
\end{ps}
\item[\code{bad header info in ...eosFreeEOS\_data...}]
  Objaw błędu z punktu~3. Usuń pliki cache zapisane przez wadliwy program i zbuduj \code{star.exe}
  ponownie (\verb!./clean!, \verb!./mk!):
\begin{ps}
Remove-Item C:\MESA\mesa-24.08.1\data\*\cache\*.bin
\end{ps}
\item[\code{star.exe} kończy się natychmiast, bez komunikatów (kod \code{0xC00000FD})]
  Przepełnienie stosu (punkt~4): ponownie \code{install\_mesa.ps1}, potem \verb!./clean!,
  \verb!./mk!.
\item[Pierwsze uruchomienie jest bardzo wolne]
  Tworzenie plików cache (jednorazowo). Rozpakowywanie i kompilację spowalnia też antywirus.
  Można (jako administrator) wyłączyć skanowanie katalogu \verb!C:\MESA!.
\end{description}

\subsection{Zawartość katalogu \texttt{windows}}

\begin{center}\small
\begin{tabular}{@{}l>{\raggedright\arraybackslash}p{0.55\textwidth}@{}}
\toprule
\textbf{Plik} & \textbf{Rola / odpowiednik na Linuksie} \\
\midrule
\code{mesa\_env.ps1} & \code{export MESA\_DIR=...}, \code{source mesasdk\_init.sh} \\
\code{install\_mesa.ps1} & \code{./install} \\
\code{test\_mesa.ps1} & testy modułów (\code{build\_and\_test}) \\
\code{makefile\_header} & \code{utils/makefile\_header} dla Windows \\
\code{forum\_makefile} & kompilacja biblioteki forum \\
\code{tools\textbackslash makedepf90.ps1} & \code{makedepf90} z MESA SDK \\
\code{tools\textbackslash cp\_if\_newer.ps1} & \code{utils/cp\_if\_newer} \\
\code{tools\textbackslash mesa\_win\_locale.c} & poprawka locale (punkt~3) \\
\code{..\textbackslash mk.ps1, rn.ps1, re.ps1, clean.ps1} & \code{./mk}, \code{./rn}, \code{./re}, \code{./clean} \\
\bottomrule
\end{tabular}
\end{center}

\end{document}
